\documentclass[a4paper, 12pt]{article}
\usepackage{wrapfig} % Для обтекания картинки текстом
\usepackage[labelformat=simple]{subcaption}
\usepackage{subcaption}
\usepackage{graphicx}
\graphicspath{{images/}} 
\usepackage{color}
\definecolor{darkblue}{rgb}{0,0,0.75}
\definecolor{darkgreen}{rgb}{0,0.75,0}
\usepackage[%pagebackref,
colorlinks, linkcolor=darkblue,
citecolor=darkgreen, urlcolor=blue, bookmarks=false, breaklinks]{hyperref}
\usepackage[utf8]{inputenc}
\usepackage[T2A]{fontenc}
\usepackage[english,russian]{babel}
\usepackage{indentfirst}
\usepackage[left=2.5cm, right=1.5cm, top=2cm, bottom=2cm]{geometry}
\usepackage{amsmath}







  \renewcommand{\contentsname}%
    {Содержание}%



\begin{document}

{\fontsize{14pt}{18}\selectfont
  
  \thispagestyle{empty}
  \begin{center}
		
		\vfill\vfill \ \\ {%\Large
			МОСКОВСКИЙ ГОСУДАРСТВЕННЫЙ УНИВЕРСИТЕТ \\
			имени М.В.~ЛОМОНОСОВА
			
			\medskip
			
			МЕХАНИКО-МАТЕМАТИЧЕСКИЙ ФАКУЛЬТЕТ


		}



{%\Large
	\vfill\vfill\vfill\vfill\vfill\vfill\vfill {%\Large
   
		Отчёт о задаче №5
		
	}
		
		
		\vfill{\Large
        
			\textbf{Алгоритм Брезенхема для рисования отрезков и окружностей.} 
		}
		
			

			\vfill
			\begin{flushright}
				\begin{tabular}{r}
                \\
                \\
                \\
                \\
                \\
                \\
                        Научный руководитель:\\
н.с. Романов Александр Вячеславович\\
                \\
					Выполнила:\\
студентка 223 группы \\
Якшибаева Дарья Сергеевна
					
				\end{tabular}
			\end{flushright}
			
		}
		
		\vfill\vfill\vfill\vfill 
		
		Москва
		
		 2025
	\end{center}
}

    \newpage
    \tableofcontents
    \fontsize{14pt}{20pt}\selectfont
    \newpage
\pagestyle{plain}

\section{Описание задачи}
 Задача состоит в написании программы на основе алгоритма Брезенхема и его аналога на языке Python.Результатом работы является графическое изображение с построенными фигурами в системе координат.\\
Алгоритм Брезенхема - алгоритм, определяющий, какие точки двумерного растра нужно закрасить, чтобы получить близкое приближение к прямой линии.
%\newpage
\section{Постановка задачи}

Входные данные:
\begin{itemize}
\setlength{\itemindent}{2cm}
\itemПараметры отрезка (координаты начала и конца)

\itemПараметры окружности (координаты центра и радиус)
\end{itemize}

Выходные данные:\\
\\
\indentГрафический файл (output.png) с изображением:
\begin{itemize}
\setlength{\itemindent}{1cm}

\itemОтрезка, построенного по алгоритму Брезенхема

\itemОкружности, построенной по алгоритму Брезенхема
\end{itemize}





\newpage
\section{Реализация задачи}

Основная идея алгоритма заключается в выборе пикселя, который находится ближе к идеальной математической линии, используя только целочисленные операции (сравнения и сложения).\\
\\
\textbf{Рассмотрим алгоритм Брезенхема для отрезка}.\\
На каждом шаге выбор состоит из двух соседних пикселей: \\
(x+1,y) - горизонтальный шаг\\
(х+1,y±1) - диагональный шаг\\
Для упрощения алгоритма будем считать, что начальная точка ($x_0$,$y_0$) лежит левее по оси X, чем конечная ($x_1$,$y_1$).В противном случае, меняем их местами:
\begin{center}
    $x_0$, $x_1$ = $x_1$, $x_0$ \\
 $y_0$, $y_1$ = $y_1$, $y_0$\\
\end{center}
 
Для обеспечения плавности и точности рисования необходимо в начале проверить угол наклона прямой. Если угол наклона между прямой и осью X больше $45^0$, то алгоритм меняет местами x и y , чтобы вести расчеты относительно оси Y, а не X.
\begin{center}
                 
         if dy>dx:\\
        $x_0$, $y_0$ = $y_0$, $x_0$\\
        $x_1$, $y_1$ = $y_1$, $x_1$\\
        dx, dy = dy, dx ,
\end{center}
гдe dy=($y_1$-$y_0$) и dx=($x_1$-$x_0$)\\
\textbf{Нахождение "ошибки":} \\
Введем вспомогательный параметр delta, измеряющий как текущее положение пикселя отклоняется от идеальной прямой.\\
Рассмотрим начальное значение delta и вычислим его.\\
Для удобства будем считать, что $y_0<y_1$\\
\[F=\frac{(x-x_0)}{(x_1-x_0)}-\frac{(y-y_0)}{(y_1-y_0)} \]
\\
\[F=y(x_1-x_0)-x(y_1-y_0)+C=ydx-xdy+C (1)\]
\newpage
Точка M(x+1,y+$\frac{1}{2}$) является серединой между точками (x+1,y) и (x+1,y+1).Подставим ее в уравнение (1).\\

\[F=dy(x+1)-dx(y+\frac{1}{2})+C\]
\\
Можем домножить на два и убрать константу, так как данные операции не повлияют на знак выражения.\\
\\
\[delta=2dy(x+1)-dx(2y+1)\]
\\
Для начальных условий:\\
\\
\[delta=2dy-dx\]
\\
Если delta>0, то математическая линия располагается выше точки M по оси Y. Следовательно выбираем пиксель (x+1,y+1).\\
Если delta<0, то математическая линия располагается ниже точки M по оси Y. Следовательно выбираем пиксель (x+1,y).\\
После выбора между горизонтальным и диагональным шагом необходимо обновить ошибку.\\
\[\Delta delta=delta_{\text{new}}-delta=2(F_{\text{new}}-F) ,\]
\textbf{Горизонтальный шаг}:\\
После горизонтального шага координата точки - (x+1,y)\\
\[\Delta delta=2(dy(x+1)-dx*(y)-dy*(x)+dx*(y))\]
\[\Delta delta=2dy\]
\[delta_{\text{new}}=delta+2dy\]\\
\newpage
\noindent\textbf{Диагональный шаг}:\\
После диагонального шага координата точки - (x+1,y+1)\\
\[\Delta delta=2(dy(x+1)-dx(y+1)-dy*(x)+dx*(y))\]
\[\Delta delta=2dy-2dx\]
\[delta_{\text{new}}=delta+2dy-2dx\]\\
Данным алгоритмом необходимо пройти от $x_0$ до $x_1$.\\
for x in range(x0, x1 + 1):\\
 if delta >= 0:\\
        \indent    delta -= 2 * dx\\
        delta += 2 * dy\\
        \\
\textbf{Рассмотрим аналог алгоритма Брезенхема для окружности.}\\
Отличие алгоритма Брезенхема для окружности от алгоритма для прямой заключается в использовании
симметрии, чтобы нарисовать только $\frac{1}{8}$ часть окружности (октант), а затем зеркально отражает эти точки для получения полной окружности.
  \[points = [
                (x_0 + x , y_0 + y), (x_0 - x , y_0 + y),\]
               \[ (x_0 + x , y_0 - y), (x_0 - x , y_0 - y),\]
                \[(x_0 + y, y_0 + x ), (x_0 - y, y_0 + x ),\]
              \[  (x_0 + y, y_0 - x ), (x_0 - y, y_0 - x )
            ]\]\\
Начальная точка:(0,r), где r - это радиус окружности.\\
Aлгоритм работает в первом октанте (от $0^0$ до $45^0$), пока x не превысит y. После этого дальнейшие точки просто зеркально повторяются.
\begin{center}
     while x <= y:\\
        plot\_points(x, y)
\end{center}

\newpage
\begin{figure}[h]
\centering
\begin{subfigure}{0.6\textwidth}
    \includegraphics[width=\linewidth]{circle.jpg}
    
\end{subfigure}
\caption{Деление окр-ти на октанты}
\end{figure}

Аналогично случаю с прямой, введём вспомогательный параметр, который измеряет как текущее положение пикселя отклоняется от математической окружности.\\
\[F=x^2+y^2-r^2(1) \]
\\
Точка M(x+1,y-$\frac{1}{2}$) является серединой между точками (x+1,y) и (x+1,y-1).Подставим ее в уравнение (1) и применим начальные условия.\\
\[F(1,r-\frac{1}{2})=1^2+(r-\frac{1}{2})^2-r^2 \]
\[F(1,r-\frac{1}{2})=-r+\frac{5}{4}\]
\[delta=5-4r\]\\
Разделим на два, так как данная операция не повлияет на знак параметра. В итоге:
\[delta=3-2r\]

Если delta>0, то точка располагается вне окружности. Следовательно выбираем пиксель (x+1,y-1).\\
Если delta<0, то точка располагается внутри окружности. Следовательно выбираем пиксель (x+1,y).\\
После выбора между горизонтальным и диагональным шагом необходимо обновить ошибку.\\
\[\Delta delta=delta_{\text{new}}-delta=\frac{4(F_{\text{new}}-F)}{2} ,\]
\textbf{Горизонтальный шаг}:\\
После горизонтального шага (x+1,y) новая средняя точка - (x+2,y-$\frac{1}{2}$)\\
\[\Delta delta=2((x+2)^2+(y-\frac{1}{2})^2-r^2-(x+1)^2-(y-\frac{1}{2})^2+r^2)\]
\[\Delta delta=2(2x+3)\]
\[delta_{\text{new}}=delta+4x+6\]\\

\noindent\textbf{Диагональный шаг}:\\
После диагонального шага (x+1,y+1) новая серединная точка - (x+2,y-$\frac{3}{2}$)\\
\[\Delta delta=2((x+2)^2+(y-\frac{3}{2})^2-r^2-(x+1)^2-(y-\frac{1}{2})^2+r^2)\]
\[\Delta delta=2(2x-2y+5)\]
\[delta_{\text{new}}=delta+4x-4y+10\]\\

\textbf{Почему разный подход к перерасчету ошибки у окружности и прямой?}\\
Прямая — это линейная функция, и ошибка меняется предсказуемо,поэтому достаточно сравнить два соседних пикселя.\\
Окружность — это квадратичная функция 
\[(x^2 + y^2 = r^2),\] и ошибка меняется нелинейно.
Если просто сравнивать два шага, можно пропустить оптимальный пиксель.
\\
\\
Для обеспечения наглядности и корректности визуализации геометрических фигур в программе реализована декартова система координат.
\\

\newpage
\section{Примеры выполнения}
\begin{minipage}{\linewidth} 
\begin{wrapfigure}{r}{0.7\textwidth} 
    \centering
    \begin{minipage}{\linewidth} 
        \includegraphics[width=\linewidth]{output.png}
        \caption{Пример №1}
        \label{fig:example1}
        
        \vspace{0.5cm} 
        
        \includegraphics[width=\linewidth]{output1.png}
        \caption{Пример №2}
        \label{fig:example2}
    \end{minipage}
\end{wrapfigure}

Начальные условия:\\
\\
$x_0$: 0\\
$y_0$: 0\\
$x_1$: 0\\
$y_1$: 200\\
\\
center X: 0\\
center Y: 0\\
radius: 200\\
\\
\\
\\
\\
\\
\\
\\
\\
$x_0$: 100\\
$y_0$: 10\\
$x_1$: -200\\
$y_1$: -100\\
\\
center X: -100\\
center Y: 50\\
radius: 100\\


\end{minipage}

\newpage
\section{Вывод}
Реализованные алгоритмы Брезенхема успешно выполняют свою задачу по визуализации геометрических фигур, используя только целочисленные операции, а именно:\\
Для отрезков:
\begin{itemize}
\itemТочно аппроксимирует прямые линии пикселями.

\itemСохраняет "гладкость" при любых углах наклона.
\end{itemize}

\noindentДля окружностей:
\begin{itemize}
\itemОбеспечивает 8-стороннюю симметрию.
\item Сохраняет "гладкость" в любом октанте.
\end{itemize}

\noindent Инструменты проверки точности алгоритмов:
 \begin{itemize}
     \itemКоординатная сетка с шагом 50пикселей
     \item Цифровые подписи осей
     \item Контрастные цвета фигур 
 \end{itemize}

\end{document}